\chapter[Classificadors de subobjectes i topos]{\texorpdfstring{Classificadors de subobjectes i introducció als topos}{Classificadors de subobjectes i introducció als topos}}

Aquest test recull els conceptes essencials del capítol. Cada pregunta té una única resposta correcta.

\begin{Form}
\iniciTest{top}{c,b,d,b,c,d,b,c,d,b}

\begin{enumerate}

\pregunta A $\cat{Set}$, la funció característica d'un subconjunt $U\subseteq E$ pren valors a:
\begin{enumerate}
  \opcio{a}{$E$.}
  \opcio{b}{$U$.}
  \opcio{c}{el conjunt $\{0,1\}$.}
  \opcio{d}{els nombres naturals.}
\end{enumerate}
\feedback

\pregunta Un \textbf{classificador de subobjectes} és un objecte $\Omega$ amb un morfisme:
\begin{enumerate}
  \opcio{a}{$\Omega\to 1$.}
  \opcio{b}{$\mathsf{true}\colon 1\to\Omega$ amb una propietat universal de pullback.}
  \opcio{c}{$\Omega\to\Omega\times\Omega$.}
  \opcio{d}{$0\to\Omega$.}
\end{enumerate}
\feedback

\pregunta La propietat universal del classificador diu que cada subobjecte $U\rightarrowtail E$ determina:
\begin{enumerate}
  \opcio{a}{un isomorfisme $E\cong\Omega$.}
  \opcio{b}{cap morfisme.}
  \opcio{c}{diversos morfismes $E\to\Omega$.}
  \opcio{d}{un únic morfisme $\chi\colon E\to\Omega$ que fa un quadrat cartesià amb $\mathsf{true}$.}
\end{enumerate}
\feedback

\pregunta L'existència d'un classificador de subobjectes equival a dir que el functor $\Sub(-)$ és:
\begin{enumerate}
  \opcio{a}{constant.}
  \opcio{b}{representable, amb objecte representant $\Omega$.}
  \opcio{c}{no functorial.}
  \opcio{d}{un isomorfisme.}
\end{enumerate}
\feedback

\pregunta El classificador de subobjectes, quan existeix, és únic:
\begin{enumerate}
  \opcio{a}{només a $\cat{Set}$.}
  \opcio{b}{mai.}
  \opcio{c}{llevat d'isomorfisme (per Yoneda).}
  \opcio{d}{llevat d'homotopia.}
\end{enumerate}
\feedback

\pregunta Un \textbf{topos elemental} és una categoria amb límits finits, exponencials i:
\begin{enumerate}
  \opcio{a}{tots els colímits infinits.}
  \opcio{b}{un objecte inicial estricte.}
  \opcio{c}{una estructura de grup.}
  \opcio{d}{un classificador de subobjectes.}
\end{enumerate}
\feedback

\pregunta Quina d'aquestes categories és un topos?
\begin{enumerate}
  \opcio{a}{la categoria dels grups.}
  \opcio{b}{$\cat{Set}$.}
  \opcio{c}{$\cat{Top}$.}
  \opcio{d}{un monoide no trivial vist com a categoria.}
\end{enumerate}
\feedback

\pregunta Per a $\cat{C}$ petita, la categoria de prefeixos $\cat{Set}^{\cat{C}\op}$ és:
\begin{enumerate}
  \opcio{a}{mai un topos.}
  \opcio{b}{un topos només si $\cat{C}$ és un grup.}
  \opcio{c}{sempre un topos.}
  \opcio{d}{cartesiana tancada però sense classificador.}
\end{enumerate}
\feedback

\pregunta A la categoria de prefeixos, el classificador $\Omega$ es construeix amb:
\begin{enumerate}
  \opcio{a}{els objectes terminals.}
  \opcio{b}{els coproductes.}
  \opcio{c}{els productes fibrats.}
  \opcio{d}{els \emph{sedassos} sobre cada objecte.}
\end{enumerate}
\feedback

\pregunta En un topos, l'exponencial $\Omega^{E}$ s'anomena:
\begin{enumerate}
  \opcio{a}{objecte inicial.}
  \opcio{b}{objecte de parts $\mathcal{P}(E)$, i classifica subobjectes de $A\times E$.}
  \opcio{c}{objecte terminal.}
  \opcio{d}{nucli de $E$.}
\end{enumerate}
\feedback

\end{enumerate}
\clauestatica
\finalitzaTest
\end{Form}
